The \asimemg is the internal instruction memory of the \aschip.

\paragraph{General Features: }
\begin{itemize}
  \item 32\;KiB (8192 words of 32\,bit), word addressed (\texttt{addr\_i[14:2]})
  \item One single-port RAM: on silicon the XH018 low-power SPRAM macro
    (Section~\ref{subsec:sramSelectionA}), in simulation a behavioural
    model with the same clocking
  \item Written only by the JTAG loader (instruction \texttt{I\_MEM});
    read-only for the CPU
  \item Synchronous read: the address is taken at the rising clock edge,
    the instruction is valid after it
  \item No reset of the contents; after power-up the contents are undefined
    until the program is loaded
\end{itemize}

\paragraph{Clocking.}
The macro has one clock input. A clock multiplexer in \texttt{as\_imem}
selects the clock with \texttt{im\_mode\_i} (from \texttt{ir\_decode},
1 while the JTAG IR holds \texttt{I\_MEM}):

\begin{table}[H]
\caption{\asimemg Operating Modes}
\label{tab:asImemModes}
\centering
\begin{tabularx}{\textwidth}{|l |l |l |X|}
  \hline
  \texttt{im\_mode\_i} & Macro clock & Access & Notes \\
  \hline
  \hline
  0 (function) & \texttt{clk\_i} (\texttt{clk\_div\_s}) & read in every cycle &
    \texttt{CEn}\,=\,0, \texttt{WEn}\,=\,1; CPU instruction fetch \\
  \hline
  1 (load) & \texttt{tck\_i} & write when \texttt{wr\_i}\,=\,1, otherwise none &
    \texttt{CEn}\,=\,\texttt{WEn}\,=\,\textasciitilde\texttt{wr\_i};
    the output keeps the last read word \\
  \hline
\end{tabularx}
\end{table}

\paragraph{Load procedure.}
The select changes when the JTAG IR is updated, asynchronously to the core
clock. The multiplexer is therefore glitch-free only if the core clock stands
still while the IR changes:
\begin{enumerate}
  \item Set \texttt{rst\_i}\,=\,1. The CGU counter is reset, so
    \texttt{clk\_core} is constant.
  \item Load the IR with \texttt{I\_MEM}, shift and update one data register
    per word (address, data, write bit).
  \item Load another instruction into the IR (e.g.\ \texttt{BYPASS}).
  \item Release \texttt{rst\_i}; the CPU starts at the reset vector.
\end{enumerate}
The first access after power-up must respect the power-up time of the macro
data sheet; with the reset sequence above this is always met.
